\scanpage{194}
\mspage{68}
% archive sheet number: 193

\deleted{Par réc\unsure{.} sur $\mathrm{card}\,A$, on est ramené au cas de \uline{deux} parties $I_1,I_2$. Supposons que $I_1$ n'est pas cofinal, prouvons que $I_2$ est cofinal. Soit $i_0\in I$, \deleted{\unsure{tel que}} tel qu'il n'existe pas d'élément de $I_1$ $\geq i_0$. On a donc $I_{/i_0}\subset I\setminus I_1$, donc $I_{/i_0}\subset I_2$ puisque $I_1\cup I_2=I$, donc $I_2$ est cofinal, $I_{/i_0}$ l'étant. \uline{NB} Cet argument montre $=$ mieux, savoir qu'il existe un $\alpha$ \inserted{et un $i_0\in I$} tel que $I_\alpha\supset I_{/i_0}$.}

Dém par l'absurde. Si pour \deleted{tt $\alpha$,} aucun $\alpha$, $I_\alpha$ \unsure{n'était|n'étant} cofinal, on en conclurait que pour tt $\alpha$ il existe un $i_\alpha$, tel que $I_{/i_\alpha}\cap I_\alpha=\varnothing$. Soit $i$ \deleted{tel que} majorant les $i_\alpha$, on aura donc $I_{/\deleted{\unsure{\alpha|i_\alpha}}i}\cap I_\alpha=\varnothing\quad\forall\alpha\in A$, contrairement à l'hypothèse que $\bigcup I_\alpha=I$.

Le cas d'une \uline{intersection} \deleted{\unsure{finie de p} \unsure{résulte|arbitraire} du lemme la définition b) \unsure{sous la} forme} de parties $L$-$\unsure{\pi|\Pi}$-accessibles : elle est bien sûr $L$-stable, donc il reste à \unsure{montrer|prouver} qu'elle est $\pi$-accessible, i.e. satisfait la condition $b_\pi$ de loc. cit.

